\section{The two different level filtrations}\label{sec:levels}
We distinguish a filtration by the number of local branches from one by the
complex dimension of a moduli stratum.  They have different indexing
categories.  Throughout this section a $q$-connected map has homotopy fibres
that are $(q-1)$-connected.  A $k$-equivalence means an equivalence on
Postnikov $k$-truncations, including components.  A colimit of fundamental
groupoids means a bicategorical colimit, with its comparison isomorphisms;
it never means the strict colimit of arbitrarily chosen one-object groups.

\subsection{Branch depth: the omitted triangle}

\begin{proposition}[The hexagon really is a counterexample]\label{P5a1}
\status{Refuted.}
For the complex manifold pair
\[
 X=\C^3,\qquad \Theta=\{z_1z_2z_3=0\},\qquad U=(\C^*)^3,
\]
the bicategorical colimit of the span in (246), cote 147, \S11,
folio 75, scan 152, has automorphism group $\Z^4$ on its single connected
component, whereas $\pi_1(U_\infty)=\Z^3$.  Thus (246), as a statement about
that span alone, is false.  The additional flags involving depth three kill
exactly the extra generator.
\end{proposition}
\begin{proof}
At depth one there are three open faces, at depth two three, and the
incidences between them form a hexagon.  Each face direction bundle has
homotopy type $T^3=(S^1)^3$.  Using the three ambient coordinate arguments
identifies all its meridians with the same ordered basis of $\Z^3$; every
incidence map is the identity on this basis.  The homotopy colimit of the
six vertex spaces and six mapping cylinders is therefore
$T^3\times S^1$.  Equivalently, choose five edges of the hexagon as a maximal
tree.  The graph-of-groups presentation identifies all vertex groups and
has one remaining edge letter $t$, with
\[
 \langle a,b,c,t\mid [a,b]=[a,c]=[b,c]
 =[t,a]=[t,b]=[t,c]=1\rangle=\Z^4.
\]
This proves the assertion about the bicategorical colimit by van Kampen.
The real oriented blow-up is
$[0,\infty)^3\times T^3$.  Its boundary has contractible radial factor
$\partial[0,\infty)^3$, so its fundamental group is $\Z^3$.
Adding the depth-three face cones the hexagon to the barycentric
subdivision of a two-simplex.  There are seven vertices, twelve edges,
and six two-simplices; the six triangular relations kill the six
non-tree edge generators in that subdivision.  The program
\texttt{code/levels\_colimit.py} builds these two order complexes,
constructs their spanning-tree presentations, and performs exact
Tietze eliminations.  Its output is respectively one and zero surviving
index-space generators, hence $\Z^4$ and $\Z^3$ after including the torus.

A strict colimit of six isomorphic one-object groups would give $\Z^3$;
that calculation discards the comparison paths around the hexagon and
is not the descent or van Kampen colimit.  Conversely, if one includes
additional two-isomorphisms from depth three as extra data, one has
already enlarged the displayed span.  The author's marginal request to
make the ``2-fl\`eches'' explicit points precisely to the missing
coherence, without supplying it in (246).
\end{proof}

\begin{theorem}[Sharp branch-depth bound]\label{P5a2}
\status{Proved.}
Let $(X,\Theta)$ be a paracompact complex normal-crossings manifold pair,
with self-intersections allowed, or a smooth analytic Deligne--Mumford
stack with this structure.  Write $M=X^{\log}$ and $L=\partial M$, using
classifying spaces in the stack case.  Let $L_{\le r}$ be the locus of
branch depth at most $r$.  For $k\ge0$ the map
\[
 L_{\le k+2}\longrightarrow L
\]
is $(k+1)$-connected and in particular a $k$-equivalence.  The homotopy
colimit of the direction bundles and all their coherent flag maps in
depths $1,\ldots,k+2$ has the same $k$-type as $L$.
The bound $k+2$ cannot be replaced uniformly by $k+1$.

This is a theorem about the spaces of the pieces.  Their fundamental
groupoids alone can replace them in this assertion if every open stratum
is aspherical; no assertion that arbitrary $k$-types follow from
fundamental groupoids is implicit.
\end{theorem}
\begin{proof}
In coordinates $z_i=r_i e^{i\theta_i}$, a depth-$d$ face is given by
$d$ radial equations $r_i=0$.  In the topological manifold
$\partial[0,\infty)^d\times\R^{2\dim_\C X-d}$ its codimension is
$d-1$, not $2d$.  The deleted locus $L\setminus L_{\le k+2}$ thus has
codimension at least $k+2$.  Triangulate the corner stratification locally.
Relative general position moves a map of a disk of dimension at most
$k+1$, whose boundary avoids that locus, off every deleted stratum:
for a source simplex of dimension $j$ and a deleted stratum of
codimension $c$, the transverse intersection has dimension at most
$j-c<0$.  The perturbations are made relative to the boundary and
successively over a locally finite subdivision.  A compact disk uses
only finitely many of the locally finite charts.  This proves
$\pi_j(L,L_{\le k+2})=0$ for $j\le k+1$, component by component, with
the analogous path-lifting assertion in degree zero.  The long exact
sequence gives the stated connectivity and equivalence of $k$-types.

For a smooth analytic Deligne--Mumford stack choose an \'etale
atlas and its \v Cech nerve $X_\bullet$.  Every $X_j$ is a manifold
with the pulled-back normal-crossings structure.  The depth-complement
inclusion on $\partial X_j^{\log}$ is $(k+1)$-connected by the
manifold argument.  Realization of a diagram of $(k+1)$-connected maps
preserves that connectivity.  One proof replaces the diagrams by
relative CW diagrams: every relative cell starts in dimension at
least $k+2$, and realization multiplies it by a simplex, only increasing
its dimension.  Equivalently, connected maps of a fixed degree are
closed under colimits in spaces.  Applying this to the \'etale nerve
gives the claimed statement on stack classifying spaces.  It does not
require finite isotropy or a proper orbifold presentation.  Branch
permutations are retained by the nerve and do not change the local
codimension calculation.

Here is the comparison with the flag diagram.  In the truncated
manifold with corners, cover the boundary by small collars of its
normalized closed faces.  An intersection indexed by a flag retracts
to the corresponding truncated face.  A truncated face is homotopy
equivalent to its open face: an inward collar vector field makes all
its remaining radial coordinates positive.  The resulting
homotopy-coherent diagram of open faces has homotopy colimit
$L_{\le k+2}$ by the collar cover.  This argument retains the normalized
branch flags, so also handles self-intersections and finite chart
symmetries.  The open face is the quotient of the punctured normal
bundle by the contractible positive radial factors.

For sharpness put $m=k+2$ and take the coordinate divisor in $\C^m$.
Keeping depths at most $m-1$ removes the radial origin from
$\partial[0,\infty)^m$.  Radial normalization $\sum r_i=1$ identifies
its homotopy type with $\partial\Delta^{m-1}=S^{m-2}=S^k$.  Thus
\[
 L_{\le k+1}\simeq T^{k+2}\times S^k,
 \qquad L\simeq T^{k+2}.
\]
For $k=0$ their component sets differ; for $k=1$ their fundamental
groups differ; for $k\ge2$ the sphere contributes a nonzero
$\pi_k$ killed by the inclusion.  The code checks the boundary-simplex
chain matrices for $0\le k\le4$; the radial argument proves every $k$.
Finally a torus bundle over an aspherical open stratum is aspherical
by its homotopy exact sequence.  Finite flag covers preserve
asphericity.  On these pieces $B\Pi_1$ is therefore a homotopy
equivalence, naturally in the homotopy-coherent sense, and can be
substituted in the homotopy colimit.
\end{proof}

\begin{corollary}[Corrected Corollaires 2 and 3]\label{P5a3}
\status{Proved.}
In the generality of Theorem~\ref{P5a2}, the fundamental groupoid of
$U_\infty$ is the bicategorical colimit of the complete coherent flag
diagram in branch depths $1,2,3$.  If no three branches meet, the span
in (246) suffices.  If $U_\infty\to U$ is a $1$-equivalence, that same
corrected colimit is $\Pi_1(U)=\Pi_1(X^*)$, not generally $\Pi_1(X)$.
\end{corollary}
\begin{proof}
Apply Theorem~\ref{P5a2} with $k=1$ and the groupoid version of van
Kampen to its collar diagram.  Fundamental groupoids preserve
homotopy colimits after passage to the $1$-type, independently of
asphericity: a functor to any groupoid is exactly compatible local
functors with the comparison isomorphisms and their cocycle conditions.
When depths three and higher are absent only the displayed span
remains.  The additional assumption identifies its answer with
$\Pi_1(U)$ by composition.
For the error in (247), take $X=\C$, $\Theta=\{0\}$.  Its link
$S^1\to\C^*$ is a homotopy equivalence, but $\pi_1(X)=0$ and the
colimit is $\Z$.  Thus the unstarred conclusion is false even when
there are no multiple crossings.  On inspecting scan 152 we read
an unstarred $X$ in (247); the mathematical correction is an added
star, not a claim that an invisible star has been deciphered.  The
struck draft on scan 151 includes depths through three and is
mathematically closer to the corrected formulation.
\end{proof}

\subsection{Dimension: the infinite multicurve poset}

For a stable oriented surface type $(g,n)$, $2g-2+n>0$, let
$d=3g-3+n$.  The $n$ markings are individually labelled punctures.
Let $\Gamma_{g,n}$ denote the pure mapping class group and
$\mathcal T_{g,n}$ Teichm\"uller space.  A multicurve is a finite set
of isotopy classes of pairwise disjoint, pairwise nonisotopic,
essential nonperipheral simple closed curves.  We allow the empty
multicurve and put
\[
 \mathcal P_k(g,n)=\{\sigma: |\sigma|\ge d-k\},
\]
ordered by inclusion.  If $k\ge d$, the empty multicurve is a minimum.
This is the infinite multicurve poset, not the finite poset of
combinatorial stable-tree types.  Even on a four-punctured sphere
there are infinitely many curve isotopy classes.

\begin{theorem}[The homotopy fibre of a dimension level]\label{P5b1}
\status{Proved.}
For every stable $(g,n)$ and integer $0\le k\le d$, choose the
standard sufficiently thin plumbing neighbourhood of the union of
Deligne--Mumford strata of dimension at most $k$, and intersect it
with $\mathcal M_{g,n}$.  Denote the resulting orbifold by $W_k$.
There is a natural weak equivalence, up to the choice of marking,
\[
 \operatorname{hofib}\bigl(BW_k\longrightarrow B\mathcal M_{g,n}\bigr)
 \simeq |\mathcal P_k(g,n)|.
\]
For $k=d$ we take $W_d=\mathcal M_{g,n}$; both sides are contractible.
Replacing this poset by its set of $\Gamma_{g,n}$-orbits does not
compute the homotopy fibre.
\end{theorem}
\begin{proof}
Fix a collar-lemma constant $\epsilon>0$ such that two geodesics
of length less than $\epsilon$ cannot intersect.  For a multicurve
$\sigma$ let
\[
 V_\sigma=\{t\in\mathcal T_{g,n}:
        \ell_\alpha(t)<\epsilon\text{ for all }\alpha\in\sigma\}.
\]
Extend $\sigma$ to a pants decomposition.  Fenchel--Nielsen
coordinates identify $V_\sigma$ with a product of intervals
$(0,\epsilon)$, intervals $(0,\infty)$, and copies of $\R$;
it is contractible.  An intersection of finitely many such sets
is $V_{\cup\sigma_i}$ if the union is a multicurve and is empty
otherwise.  These coordinate facts, including the description of
nodes by vanishing lengths, are stated in \citet[Chapter 1, \S2, Theorem 1.1]{Harer1988} and the paragraph following it; we have applied
them only after extending a specified multicurve to a pants
decomposition.

Let $V_k=\bigcup_{|\sigma|\ge d-k}V_\sigma$ when $k<d$.
The homotopy colimit of this system maps to $V_k$ by evaluation.
Over a point $t$ its incidence category consists of subsets, of
size at least $d-k$, of the finite set of all $\epsilon$-short
curves of $t$.  It has that full short-curve set as a maximum,
so its nerve is contractible.  Equivalently the good-cover nerve
lemma, followed by the subdivision indexed by unions, proves
\[
 V_k\simeq |\mathcal P_k(g,n)|.
\]
The construction is $\Gamma_{g,n}$-equivariant.  The plumbing
interpretation of Fenchel--Nielsen coordinates identifies
$[V_k/\Gamma_{g,n}]$ with a sufficiently thin model of the requested
neighbourhood: a curve with $c$ nodes has moduli dimension $d-c$,
so $c\ge d-k$ is exactly the level condition.  Shrinking the
collars changes neither the homotopy type nor the map to the
open moduli stack.  One can alternatively take this short-geodesic
model as the precise meaning of its thickening.

Teichm\"uller space is contractible, again by the same coordinates.
The map of Borel spaces
\[
 E\Gamma\times_\Gamma V_k\longrightarrow
 E\Gamma\times_\Gamma\mathcal T_{g,n}\simeq B\Gamma
\]
has homotopy fibre $V_k$.  This proves the assertion, with all
finite automorphism groups retained.  It also explains why taking
a coarse quotient before forming the fibre is invalid.  For
$k=d$ the set $V_\varnothing$ is all of Teichm\"uller space and the
poset has a minimum, proving the endpoint assertion.
\end{proof}

\begin{theorem}[Genus-zero connectivity, with the endpoint corrected]\label{P5b2}
\status{Known, with reduction proved.}
For $n\ge3$, $d=n-3$, and every $k\ge0$, the poset
$\mathcal P_k(0,n)$ is $(k-1)$-connected.  Consequently
$W_k\to\mathcal M_{0,n}$ is $k$-connected.
For $k\ge d$ the poset must include the empty multicurve; it is
then contractible.  If one excludes the empty multicurve at
$k=d$, the assertion is false, already for $(n,k)=(4,1)$ and
also for the requested test $(n,k)=(6,3)$.
\end{theorem}
\begin{proof}
Let $K=\mathcal C_{0,n}$ be the unmodified disjointness curve
complex: a simplex of dimension $s-1$ is a nonempty multicurve
with $s$ components.  In particular, $\mathcal C_{0,4}$ is an
infinite discrete set, not the Farey graph obtained by changing
the edge rule.  The Harer theorem, in the directly checked form
\citet[Chapter 4, Theorem 4.1]{Harer1988}, together with Chapter 3,
Lemma 3.6, gives $K\simeq\bigvee S^{n-4}$ for $n\ge4$.
Only the resulting $(n-5)$-connectivity is needed here.

We verify the requisite assertion for every link.  Cutting along
$s$ curves gives $s+1$ genus-zero components with $n_j$ special
points, where $\sum_j n_j=n+2s$.  A further curve is exactly a
curve in one of those components, and disjointness is componentwise.
Thus the link is the join of their curve complexes.  Components
with three special points contribute no curves and are omitted,
or equivalently contribute the augmented complex $S^{-1}$.
The connectivity formula for joins, applied to Harer's theorem
on each remaining component, gives link dimension
\[
 \sum_j(n_j-3)-1=n-s-4=d-s-1
\]
and connectivity $d-s-2$.  This proves homotopy
Cohen--Macaulayness of $K$, including the empty links of maximal
pants decompositions.  We now give the particular rank-selection
argument, rather than just invoking its name.

Put $r=d-k>0$ and begin with the barycentric subdivision of $K$.
Delete vertices corresponding to simplices of sizes
$1,2,\ldots,r-1$, in this order.  At the stage of deleting a
size-$s$ vertex, all smaller faces have gone, so its remaining
link is the barycentric subdivision of the original link of
that simplex.  It is $(d-s-2)$-connected, and in every stage
$s\le r-1$ this is at least $(k-1)$-connected.  Reattaching a
vertex is coning off that link.  A cone on an $(k-1)$-connected
CW complex can, relative to the link, be built with cells of
dimension at least $k+1$; hence deletion preserves all
connectivity assertions through degree $k-1$.
Vertices of the same rank are pairwise nonadjacent, so all can
be treated at once.  For an infinite complex the same conclusion
follows by cellular approximation of compact disks, which uses
finitely many cells, followed by the corresponding relative
homotopies.  The remaining complex is $|\mathcal P_k(0,n)|$.
Since $K$ is $(d-2)$-connected, this proves the claim for $k<d$.
This is also the instance of rank selection covered by
\citet[Corollary 4.4]{BjornerWachsWelker2009}; the argument
above specifies the rank interval and treats the infinite case.

For $k\ge d$, the minimum $\varnothing$ cones the whole order
complex.  For $n=3$ this is the entire argument.  Without that
minimum, at $(n,k)=(4,1)$ the complex is disconnected, and at
$(6,3)$ it is a noncontractible wedge of two-spheres, not
$2$-connected.  The map connectivity follows from
Theorem~\ref{P5b1} and our stated convention.
\end{proof}

\begin{proposition}[Independent low-dimensional verification]\label{P5b3}
\status{Proved.}
The cases $(0,5),k=1$ and $(0,6),k=2$ of Theorem~\ref{P5b2}
follow independently of the curve-complex and rank-selection
theorems from an explicit strictly plurisubharmonic exhaustion.
The case $(0,6),k=3$ is contractible by the empty-multicurve cone.
The exact algebraic checks are recorded in
\texttt{code/levels\_stein.py} and its output.
\end{proposition}
\begin{proof}
Fix three markings at $\infty,0,1$.  For $d=n-3$ this gives
\[
 \mathcal M_{0,n}=\{(z_1,\ldots,z_d)\in\C^d:
 z_i\ne0,1,\ z_i\ne z_j\}.
\]
Let $\mathscr A_d$ be the list of linear forms $z_i,z_i-1,z_i-z_j$
with $i<j$, and put
\[
 f(z)=\sum_i|z_i|^2+\sum_{\ell\in\mathscr A_d}|\ell(z)|^{-2}.
\]
This is proper: its first term controls escape to Euclidean
infinity and its other terms control approach to the deleted
hyperplanes.  If $a_\ell$ is the row of coefficients of $\ell$,
its Levi matrix is exactly
\[
 I_d+\sum_{\ell\in\mathscr A_d}
       |\ell|^{-4}a_\ell^*a_\ell.
\]
For every nonzero vector $v$ its quadratic value is
$\|v\|^2+\sum_\ell|\ell|^{-4}|a_\ell v|^2>0$.
The code independently expands every principal minor, treating
the positive numbers $|\ell|^{-4}$ as independent indeterminates.
For $d=2$ there are five forms and the determinant has fourteen
positive monomials.  For $d=3$ there are nine forms and the
determinant has ninety-three positive monomials.  Every
principal minor has constant term one and strictly positive
coefficients.  These are exact polynomial calculations, not
numerical samples of an infinite curve complex.

Add a generic real linear function to $f$.  Strict Levi
positivity and properness remain true, and Sard's theorem
applied to the gradient makes the perturbed function Morse.
Its critical locus is a discrete semialgebraic set, hence
finite.  At a critical point let $Q$ be its real Hessian.  Levi
positivity says $Q(v)+Q(Jv)>0$ for every nonzero $v$.  If a
negative subspace had dimension greater than $d$, it would
intersect its $J$-translate nontrivially; for a nonzero vector
in that intersection both summands would be negative.  Thus
every Morse index is at most $d$.  This is the index argument
underlying \citet[Theorem 7.2]{Milnor1963}; the handle comparison
uses Theorem 3.2 and Remark 3.3 of the same source.

Choose a regular value above every critical value.  The
sublevel is a compact core $K$, and its boundary represents the
end of this algebraic manifold: outside the core the normalized
gradient integrates to a collar because there are no further
critical values.  The end represented by a small punctured
neighbourhood in any smooth normal-crossings compactification
is the same end.  Reversing the finite Morse handle decomposition
builds $K$ from its boundary by handles of indices
$2d-\lambda\ge d$.  Therefore $\partial K\to K\simeq
\mathcal M_{0,n}$ is $(d-1)$-connected.
When $n=5,k=1$ or $n=6,k=2$, the selected strata are all boundary
strata, since $k=d-1$.  The thickening in
Theorem~\ref{P5b1} represents this end, so its fibre is
$(d-2)$-connected: connected in the first case and simply
connected in the second.  This proves exactly the requested
checks by a second method.  Finally $k=d=3$ includes the open
stratum and the empty multicurve, giving a contractible fibre.
\end{proof}

\subsection{Presentations and the meaning of two levels}

\begin{theorem}[The complex-of-groups presentation at dimension two]\label{P5c1}
\status{Proved.}
For $n\ge5$, the pure punctured-sphere mapping class group
$\Gamma_{0,n}$ is the fundamental group of the developable
complex of groups obtained from its action on
$|\mathcal P_2(0,n)|$.  Vertex groups are stabilizers of
multicurves whose moduli strata have dimension at most two;
their inertia subgroups are the free abelian multitwist groups.
A spanning tree gives the presentation described below.
Connectedness of $\mathcal P_1$ supplies generators from
strata of dimensions zero and one and their incident paths;
simply connectedness of $\mathcal P_2$ supplies the required
relations, together with the internal relations of the groups
of the pieces.  All incidence maps and twisting elements are
part of the data.
\end{theorem}
\begin{proof}
Here is a formulation avoiding an unspecified complex of
groups.  Let $Y=|\mathcal P_2(0,n)|$ and $G=\Gamma_{0,n}$.
Barycentric simplices have strictly increasing numbers of
curves, so the action has no inversions.  Choose a representative
$\widetilde v$ for each vertex of $Y/G$, representatives of
incidence edges, and group elements carrying their endpoints
to the chosen representatives.  Put $G_v=\operatorname{Stab}_G
(\widetilde v)$.  Orient an incidence edge from the larger multicurve
$\tau$ to a contained multicurve $\sigma$; a pure genus-zero mapping
class stabilizing $\tau$ fixes its curves individually and stabilizes
$\sigma$.  The chosen transport elements give injective
homomorphisms $\psi_a:G_{i(a)}\longrightarrow G_{t(a)}$ between
the corresponding stabilizers
and twisting elements $c_{a,b}$ for composable edges; by their
construction inside $G$ they satisfy the usual cocycle equation.
Fix a maximal tree $T$ in the quotient one-skeleton.
The presentation has every element of each $G_v$ and one
oriented edge letter $t_a$ as generators, with relations
\[
 \begin{gathered}
 \text{multiplication in each }G_v,\qquad
 t_{\bar a}=t_a^{-1},\qquad t_a=1\ (a\in T),\\
 t_a g t_a^{-1}=\psi_a(g),\qquad
 t_a t_b=c_{a,b}t_{ab}
 \quad\text{for each incidence triangle.}
 \end{gathered}
\]
Orientation choices can replace these equations by their
inverses; the displayed convention defines ours.  The natural
homomorphism to $G$ sends the stabilizer generators to their
actual mapping classes and the edge letters to their chosen
transport elements.

To prove surjectivity, connect a chosen lift of a vertex to its
translate by an arbitrary $g\in G$ through the connected space
$Y$.  Subdivision turns the path into an edge path, and successive
transport expresses $g$ by the displayed letters and a final
vertex stabilizer.  To prove injectivity, a word mapping to one
lifts to a loop in $Y$.  By Theorem~\ref{P5b2}, $Y$ is simply
connected.  A simplicial null-homotopy decomposes into edge
cancellations and triangle moves; these are exactly the
presentation relations.  This proves the claim without a
faithfulness assumption about a coarse quotient.  The same
argument for the connected subspace $|\mathcal P_1|$ gives generation
by its vertex groups and edge transports.  A new dimension-two vertex
group introduces no necessary new generator: its base is either
$\mathcal M_{0,5}$ or $\mathcal M_{0,4}^2$, and its boundary
thickening surjects on its fundamental group by the same theorem
applied to that two-dimensional pair (or by the product argument).
The direction-torus generators are already present over deeper
strata.  Thus its generators can be eliminated in favour of dimension
zero and one generators, while its internal relations supply the
new relations.  This proves the stated generation assertion.  Local surface cutting identifies the
stabilizer with the mapping classes of the cut components,
extended by the independent twists about the cut curves.
Independence follows by testing on arcs transverse to one cut
curve at a time; a nonzero twist exponent changes that arc's
relative isotopy class.

The quotient is finite and indexed by stable labelled trees,
but its stabilizer groups and transport elements are indispensable.
The finite tree complex alone is not the development $Y$.
For $n=5$, $\mathcal P_2$ includes the empty multicurve and its
vertex group is $G$; the explicit presentation below calculates
that group rather than using this tautological vertex as an
uncomputed input.  For $n=6$, $\mathcal P_2$ consists of the
nonempty multicurves and no such vertex occurs.
\end{proof}

We now write finite presentations with all words specified.
They are a concrete reduction of the group just reconstructed,
and fix the choices of base points and generators.  Write
\[
 A_{ij}=\sigma_{j-1}\cdots\sigma_{i+1}\sigma_i^2
             \sigma_{i+1}^{-1}\cdots\sigma_{j-1}^{-1}
 \quad(1\le i<j)
\]
in the usual disk braid group, with real ordered base
configuration.  We retain the subgroup in which the first
two strands have trivial braid, so that their positions may be
fixed at $0$ and $1$.  For words $u,v$ abbreviate
\[
 c(u,v)=v^{-1}u^{-1}vu.
\]
This is only a written-word abbreviation, not an extra generator.

\begin{proposition}[Five labelled punctures: five generators, six relations]\label{P5c2}
\status{Proved.}
For $a=A_{13}$, $b=A_{23}$, $x=A_{14}$, $y=A_{24}$,
$z=A_{34}$, a presentation of $\Gamma_{0,5}$ is
\[
 \langle a,b,x,y,z\mid\text{the following six equations}\rangle:
\]
\begin{align*}
 axa^{-1}&=z^{-1}xz,&
 aya^{-1}&=c(x,z)y c(x,z)^{-1},&
 aza^{-1}&=c(x,z)z,\\
 bxb^{-1}&=x,&
 byb^{-1}&=z^{-1}yz,&
 bzb^{-1}&=c(y,z)z.
\end{align*}
In particular $\Gamma_{0,5}=F(x,y,z)\rtimes F(a,b)$ for the
specified action.  The presentation complex is an aspherical
finite two-dimensional complex.
\end{proposition}
\begin{proof}
Forgetting the last moving point in
$\operatorname{Conf}_2(\C\setminus\{0,1\})$ is a locally trivial
fibration with fibre a plane minus three points and base a plane
minus two.  Local triviality follows by choosing pairwise
disjoint disks about the retained points and moving their
centres by ambient isotopies supported in those disks.  This
is the instance of \citet[Theorem 1]{FadellNeuwirth1962},
needed here.  It has a section: choose the new point on the
positive real axis at a radius exceeding the absolute values
of every retained point and $1$, continuously in the configuration.
Both base and fibre are aspherical graphs.  The homotopy exact
sequence gives the split extension displayed above.

For the action, use the Artin substitutions
\[
 \rho(\sigma_i):\quad
 X_i\mapsto X_iX_{i+1}X_i^{-1},\quad
 X_{i+1}\mapsto X_i,
\]
with all other $X_j$ fixed, composing functions in the usual
right-to-left order.  Positive conjugation by a standard base
braid on the kernel basis $A_{j,4}$ is the Artin substitution
for its mirror word, obtained by replacing each $\sigma_i$ by
$\sigma_i^{-1}$.  Substituting the two words $A_{13}$ and
$A_{23}$ and freely reducing gives exactly the six equations
above.  This sign convention is checked independently in
\texttt{levels\_braids.py}: each side is expanded as a braid on
four strands and the two resulting free-group automorphisms
are equal.  The Artin representation is faithful; the checked
source \citet[Theorem A, case (1), parameter $1$]{Shpilrain2001}, supplies this fact.  The semidirect-product normal form
shows that the six conjugation equations, with no relators
inside either free group, are a complete presentation.

Finally a graph bundle over a graph is built by attaching
mapping cylinders of the fibre homotopy equivalences.  Its
universal cover is a tree of contractible universal covers of
the fibre graphs, hence contractible.  Its cells are one
vertex, five edges, and six two-cells with precisely these
attaching words.  This proves the last assertion.
\end{proof}

\begin{proposition}[Six labelled punctures: nine generators, twenty-six relations]\label{P5c3}
\status{Proved.}
Use the five generators and six relations of
Proposition~\ref{P5c2}, and add
\[
 p=A_{15},\quad q=A_{25},\quad r=A_{35},\quad s=A_{45}.
\]
For each row $h$ and column $u$ in the following table impose
$huh^{-1}=$ the displayed entry.  These twenty equations,
together with the previous six, are a presentation of
$\Gamma_{0,6}$.
\[
\begin{array}{c|llll}
 &p&q&r&s\\ \hline
 a&r^{-1}pr&c(p,r)q c(p,r)^{-1}&c(p,r)r&s\\
 b&p&r^{-1}qr&c(q,r)r&s\\
 x&s^{-1}ps&c(p,s)q c(p,s)^{-1}&c(p,s)r c(p,s)^{-1}&c(p,s)s\\
 y&p&s^{-1}qs&c(q,s)r c(q,s)^{-1}&c(q,s)s\\
 z&p&q&s^{-1}rs&c(r,s)s
\end{array}
\]
Thus
$\Gamma_{0,6}=F(p,q,r,s)\rtimes
(F(x,y,z)\rtimes F(a,b))$ with the specified actions.
There is an aspherical finite three-dimensional complex with
this two-skeleton and cell counts $(1,9,26,24)$.
\end{proposition}
\begin{proof}
Apply the same forgetting fibration and section to three
moving points.  Its fibre is a plane minus four points.
The base is the aspherical space of
Proposition~\ref{P5c2}; the homotopy exact sequence therefore
makes the fibre group injective and gives the asserted split
extension.  The mirror-Artin substitutions for
$A_{13},A_{23},A_{14},A_{24},A_{34}$ on a free group of rank four
are the twenty table entries.  The code checks all of them as
actual identities of braids on five strands.  The split
extension normal form proves completeness of the presentation.

To see the cell count and the first syzygies, let $R_4$ be the
rose with four oriented edges and a chosen vertex.  Each component of
$\operatorname{Map}_*(R_4,R_4)=(\Omega R_4)^4$ is contractible, since
$R_4$ is aspherical.  Thus the monoid of based homotopy equivalences
has contractible components indexed by $\operatorname{Aut}(F_4)$.
The exact action just computed determines, up to fibre homotopy, a
graph bundle with section over the finite aspherical two-complex of
Proposition~\ref{P5c2}.  This statement can also be constructed cell
by cell: choose based graph maps on the base edges and homotopies
realizing the six exact action relations on its two-cells.  The
contractibility just proved supplies all required compatibility.

For a base $q$-cell, first attach its section cell.  Its remaining
relative fibre pair is
\[
 (D^q\times R_4,
  (S^{q-1}\times R_4)\cup(D^q\times\{*\})),
\]
whose quotient is $S^q\wedge R_4$, a wedge of four
$(q+1)$-spheres.  Consequently exactly four further cells of degree
$q+1$ are attached.  For $q=1$ their attaching words are exactly the
twenty conjugation equations in the table; for $q=2$ they are
three-cells.  This gives a homotopy CW model with the displayed
two-skeleton and cell counts
\[
 (1+4t)(1+5t+6t^2)=1+9t+26t^2+24t^3.
\]
The six base relators act identically on every fibre generator;
these are the twenty-four compatibility homotopies that attach
the three-cells.  The code verifies all twenty-four word identities
in the free group on $p,q,r,s$, and records the six base words in
\texttt{levels\_braids\_actions.json}.  Based graph maps inducing the same homomorphism
have contractible spaces of based homotopies, as follows by
lifting to the tree universal cover.  Hence these homotopies
exist coherently and no further cells are required.  The
fibration has aspherical base and fibre, so its total space is
aspherical by the homotopy exact sequence.  Its universal-cover
cellular complex has no cells above dimension three and is
acyclic.  Consequently its twenty-four three-cells give a free
$\Z[\Gamma_{0,6}]$-basis of the second homotopy module of the
displayed presentation two-complex.  This is a precise first
relations-among-relations statement for these chosen words.
\end{proof}

\begin{theorem}[The genus-zero two-level statement and its scope]\label{P5c4}
\status{Proved.}
For the pure punctured-sphere groups, gluing all strata of
modular dimension at most two, with their direction tori,
fundamental groupoids, incidence maps, and two-dimensional
coherence, determines $\Gamma_{0,n}$.  Its indecomposable
positive-dimensional building blocks at those levels are
$\mathcal M_{0,4}$ and $\mathcal M_{0,5}$; the other
dimension-two possibility is
$\mathcal M_{0,4}\times\mathcal M_{0,4}$.
At dimension three the support types are
\[
 \mathcal M_{0,6},\qquad
 \mathcal M_{0,5}\times\mathcal M_{0,4},\qquad
 (\mathcal M_{0,4})^3.
\]
They supply the next coherence needed to remove the second
homotopy left at level two.  This does not exclude syzygies
already internal to the direction tori or the pieces.

The four elementary types actually named in cote 147,
\S9, folio 29, scan 60, are
$(0,4),(1,1),(0,5),(1,2)$.  The statement above establishes
the genus-zero part of that expectation; it is not a statement
about profinite or Galois actions.
\end{theorem}
\begin{proof}
Every genus-zero boundary stratum is a product
$\prod_v\mathcal M_{0,n_v}$ with dimension
$\sum_v(n_v-3)$.  The nonnegative partitions of dimensions one,
two, and three give exactly the listed factor types, after
suppressing zero-dimensional $\mathcal M_{0,3}$ factors.
All these factors are aspherical by induction using the
forgetting fibrations in the preceding proofs; their torus
bundles are also aspherical.  Theorem~\ref{P5b2} says that the
dimension-two thickening has the same fundamental group as
$\mathcal M_{0,n}$.  Van Kampen and
Theorem~\ref{P5c1} therefore reconstruct that group from the
stated pieces and coherence.  The dimension-one thickening
surjects on fundamental groups, whereas the dimension-two
thickening accounts for all relations.  At dimension three
the inclusion is $3$-connected, so its second homotopy agrees
with that of the aspherical open moduli space and vanishes.
Thus the new compatibility data can be supplied on the three
listed support types.  At $n=6$ the only dimension-three
stratum is the open $\mathcal M_{0,6}$, and
Proposition~\ref{P5c3} makes the twenty-four syzygies explicit
as fibre-generator/base-relator pairs.

There is a complementary presentation in terms of pants moves:
the elementary moves are supported on $(0,4)$, the essential
new connected-support relation on $(0,5)$ is the pentagon,
and disjoint moves commute.  The checked \citet[\S2, Theorem 2 and Remark 4, p.~9]{HatcherLochakSchneps2000} says that the pants complex with its five specified relation types is simply connected.  Its
genus-zero restriction leaves the triangle, pentagon, and
disjoint-support square, precisely the support dimensions just
listed.  This independent formulation also makes clear why
one must include transports and coherence, not just abstract
isomorphism types of the small groups.

Finally, if ``holes'' means boundary circles fixed pointwise,
cap them to labelled punctures.  The resulting mapping class
group is a central extension of $\Gamma_{0,n}$ by the $n$
boundary twists.  For genus zero and $n\ge3$ it is noncanonically
$\Z^n\times\Gamma_{0,n}$: keep one boundary as the outside of
a disk, identify the others with framed punctures, trivialize
their tangent framings in the plane, and split the disk pure
braid group into its full-twist centre and the punctured-sphere
group.  Changing the framings changes that splitting.
The presentations above concern punctures; the boundary-fixed
version is obtained by adjoining $n$ central generators.
\end{proof}
